\section{Coleta de lixo: um heap ou um heap por processo}
\label{sec:vms-gc}

\paragraph{HotSpot.}
Um heap para todas as threads (§2.5.3 da \citetitle{doc:jvms25}) e um
coletor que cuida dele inteiro. A seção~\ref{sec:gc} mostrou as opções do
JDK~25: G1 por padrão em máquinas server-class (duas ou mais CPUs e
1792~MB ou mais; nas outras, Serial), ZGC para pausas curtas, Parallel e Serial para
casos específicos. Como o heap é um só, o coletor precisa lidar com objetos
de todas as threads ao mesmo tempo, e o tamanho do heap é uma decisão de
toda a aplicação.

\paragraph{BEAM.}
A documentação do coletor da BEAM descreve outra arquitetura:
\enquote{a per process generational semi-space copying collector using
Cheney's copy collection algorithm together with a global large object
space} (\citetitle{doc:erlang-gc}). Em português:
\begin{itemize}
\item cada processo tem a sua pilha e o seu heap, e o coletor trabalha
  sobre o heap de \emph{um} processo de cada vez;
\item dentro desse heap, a coleta é geracional e por cópia;
\item binários com mais de 64 bytes ficam fora dos heaps dos processos,
  num espaço global contado por referência; o processo guarda só um
  ponteiro para eles.
\end{itemize}
O \citetitle{livro:stenman} completa: quando um processo fica sem espaço
no heap, a VM primeiro tenta recuperar espaço coletando o lixo
\emph{daquele} processo. Como as mensagens são copiadas
(seção~\ref{sec:vms-concorrencia}), não há ponteiro de um heap de
processo para outro, e é isso que torna a coleta por processo possível.
O que a pesquisa encontrou ao testar heaps por thread na JVM, e o que a
HotSpot faz em vez disso (o TLAB), está na seção~\ref{sec:gc}.

\paragraph{O que isso muda na prática.}
O custo da coleta passa a ser proporcional ao processo, e um processo
com muita memória não obriga os outros a pagar por ele. A página do
coletor, porém, não diz que a coleta de um processo nunca afeta os
outros, e nenhuma das fontes abertas aqui promete isso; não leia
\enquote{GC por processo} como \enquote{nenhuma pausa global}. O lado
ruim aparece nos binários. O \emph{BEAM Book} explica que um pedaço
extraído de um binário (\emph{sub binary}) é só uma estrutura com
ponteiros para o binário original e aumenta a contagem de referências
dele: enquanto o pedaço pequeno estiver vivo, o binário grande também
estará.

\paragraph{GraalVM.}
No modo JIT, os coletores são os da HotSpot. No Native Image, a página
\citetitle{doc:graalvm-memory} lista outros: o Serial GC é o padrão,
descrito como \enquote{a simple (non-parallel, non-concurrent) stop and
copy GC}, com geração jovem e velha; o G1 existe só no Oracle GraalVM
(não na Community Edition), em Linux AMD64 e AArch64, com
\texttt{--gc=G1}; e o Epsilon não coleta nada. A
seção~\ref{sec:memoria-container} mostrou o efeito disso em container.

\paragraph{Ponte com o seu código.}
Em Java, o mais perto de \enquote{heap por processo} é \emph{não}
compartilhar estado mutável entre tarefas: objetos imutáveis
(\texttt{record}, seção~\ref{sec:records}) e mensagens copiadas entre
componentes. Isso não muda o GC da HotSpot, mas tira o estado
compartilhado do caminho, que é o que a BEAM resolve na VM.
